\chapter{Related Work}
\label{ch:related}

\section{Convolutional Neural Networks}
\label{sec:rw_cnn}
A CNN applies small learned filters across the whole image~\cite{lecun1998}.
A convolutional layer with $C_{\mathrm{in}}$ input channels,
$C_{\mathrm{out}}$ output channels and a $k \times k$ kernel has
$C_{\mathrm{out}}(C_{\mathrm{in}}k^2 + 1)$ parameters, independently of the
image size, and recognises a feature wherever it appears. Each convolution
is followed by a non-linearity, usually $\mathrm{ReLU}(z) = \max(0,
z)$~\cite{nair2010}, and by pooling that reduces the resolution; fully
connected layers then map the features to class scores~\cite{goodfellow2016}.
For a binary task the output $z$ is turned into a probability
$\hat p = \sigma(z) = 1/(1+e^{-z})$ and the network is trained by minimising
the binary cross-entropy
\begin{equation}
\mathcal{L} = -\frac{1}{B}\sum_{b=1}^{B}\Big[y_b \log \hat{p}_b + (1-y_b)\log(1-\hat{p}_b)\Big]
\label{eq:bce}
\end{equation}
over mini-batches of $B$ labelled examples, with gradients from
backpropagation and a stochastic optimiser such as Adam~\cite{kingma2015}.

Deep architectures such as VGG~\cite{simonyan2015} and
ResNet~\cite{he2016}, whose residual connections make very deep networks
trainable, were developed on ImageNet~\cite{deng2009}. Their early layers
learn generic features that transfer well to other tasks~\cite{yosinski2014},
so fine-tuning a pretrained network is the standard approach when labelled
medical data is scarce~\cite{litjens2017}.

\section{Deep Learning for Tuberculosis Detection}
\label{sec:rw_tb}
The Montgomery and Shenzhen datasets of Jaeger et al.~\cite{jaeger2014}, with
a few hundred radiographs each, were the first widely used TB benchmarks.
Lakhani and Sundaram~\cite{lakhani2017} fine-tuned ImageNet-pretrained
AlexNet and GoogLeNet on about a thousand radiographs and reached an area
under the ROC curve of about 0.99 with an ensemble of the two. Pasa et
al.~\cite{pasa2019} showed that a compact CNN with far fewer parameters
reaches competitive accuracy, which supports the premise of this project
that TB detection does not require a very large network.

Rahman et al.~\cite{rahman2020} assembled the dataset used here, the TB
Chest X-ray Database, from several public sources (3,500 normal and 700 TB
images). They compared nine pretrained CNNs with and without lung
segmentation and reported accuracies of about 98--99\%. Two points from
their work recur in this report: the images come from different sources
and acquisition conditions, and the best results use large networks on
images of a few hundred pixels per side.

\section{Quantum Computing Background}
\label{sec:rw_qc}
A qubit is a two-level system with basis states $\ket{0}$ and $\ket{1}$ and
general state
\begin{equation}
\ket{\psi} = \alpha\ket{0} + \beta\ket{1}
= \cos\tfrac{\theta}{2}\ket{0} + e^{i\phi}\sin\tfrac{\theta}{2}\ket{1},
\qquad |\alpha|^2 + |\beta|^2 = 1,
\end{equation}
a point on the Bloch sphere~\cite{nielsen2010}. A register of $n$ qubits has
$2^n$ complex amplitudes, which is why quantum systems are hard to simulate
and why they may represent functions that are expensive for classical
models~\cite{biamonte2017}. Computation uses unitary gates; the rotations
used here are
\begin{equation}
R_Y(\theta) = \begin{pmatrix}\cos\frac{\theta}{2} & -\sin\frac{\theta}{2}\\[2pt]
\sin\frac{\theta}{2} & \cos\frac{\theta}{2}\end{pmatrix},\qquad
R_Z(\theta) = \begin{pmatrix}e^{-i\theta/2} & 0\\ 0 & e^{i\theta/2}\end{pmatrix},
\end{equation}
\begin{equation}
\mathrm{Rot}(\phi,\theta,\omega) = R_Z(\omega)R_Y(\theta)R_Z(\phi),
\end{equation}
together with $R_X(\theta) = e^{-i\theta X/2}$, where $X, Y, Z$ are the Pauli
matrices. The two-qubit CNOT gate flips its target when the control is
$\ket{1}$; it creates \emph{entanglement}, correlations that cannot be
described qubit by qubit.

Measuring a qubit returns 0 or 1 with probabilities $|\alpha|^2$ and
$|\beta|^2$, so quantum models output \emph{expectation values} estimated
from many repeated runs (\emph{shots}):
\begin{equation}
\expval{Z_k} = \bra{\psi} Z_k \ket{\psi} = P(0) - P(1) \in [-1, 1].
\label{eq:expz}
\end{equation}
With $N$ shots the estimate has variance $(1-\expval{Z_k}^2)/N$. Real devices
are noisy and their state is a density matrix $\rho$; a common noise model
is the depolarising channel
\begin{equation}
\mathcal{E}(\rho) = (1-p)\,\rho + \frac{p}{3}\left(X\rho X + Y\rho Y + Z\rho Z\right),
\label{eq:depol}
\end{equation}
which drives every $\expval{Z_k}$ towards zero.

\section{Variational Quantum Circuits}
\label{sec:rw_vqc}
A variational circuit applies a data encoding $S(\mathbf{x})$ and a
trainable unitary $W(\boldsymbol{\theta})$ and measures,
\begin{equation}
f_k(\mathbf{x};\boldsymbol{\theta}) =
\bra{0}^{\otimes n} S(\mathbf{x})^\dagger W(\boldsymbol{\theta})^\dagger\,
Z_k\, W(\boldsymbol{\theta}) S(\mathbf{x}) \ket{0}^{\otimes n},
\label{eq:vqc}
\end{equation}
and its outputs are fed to a classical loss. Schuld et
al.~\cite{schuld2020circuit} proposed ``circuit-centric'' classifiers built
from layers of general rotations and CNOTs at increasing distances, available
in PennyLane as \texttt{StronglyEntanglingLayers}; Mitarai et
al.~\cite{mitarai2018} showed that such circuits are universal function
approximators when inputs are encoded non-linearly. Because the state
cannot be read out on hardware, gradients are obtained with the
\emph{parameter-shift rule}~\cite{mitarai2018, schuld2019grad},
\begin{equation}
\frac{\partial f}{\partial \theta_j} =
\frac{1}{2}\left[ f\!\left(\theta_j + \tfrac{\pi}{2}\right) -
f\!\left(\theta_j - \tfrac{\pi}{2}\right) \right],
\label{eq:pshift}
\end{equation}
which is exact and costs two circuit runs per parameter. In simulation,
backpropagation through the simulator gives the same gradient more cheaply.

\label{sec:rw_encoding}
The encoding largely determines what the circuit can
compute~\cite{schuld2019feature}. \emph{Angle encoding} places one feature on
each qubit, $\bigotimes_k R_Y(x_k)\ket{0}$; it needs $n$ qubits for $n$
features and one gate per qubit. \emph{Amplitude encoding} writes a
normalised vector of $2^n$ features into the amplitudes,
$\ket{\psi(\mathbf{x})} = \lVert\mathbf{x}\rVert^{-1}\sum_i x_i\ket{i}$; it is
exponentially compact but needs deep state preparation and discards the
vector's length.

\section{Hybrid Quantum--Classical Networks}
\label{sec:rw_hybrid}
Mari et al.~\cite{mari2020} formalised classical-to-quantum transfer
learning: a pretrained network is truncated to a short feature vector that
a ``dressed'' circuit, with classical layers before and after it,
classifies. Our \cq{} model follows this design. Henderson et
al.~\cite{henderson2020} introduced quanvolutional networks, in which a
random, untrained circuit transforms each image patch; on MNIST these
networks outperformed a plain CNN but performed similarly to a CNN with
random non-linear classical filters, a finding that motivates the
classical controls of our \qc{} model. The quantum convolutional network of
Cong et al.~\cite{cong2019} is a different construction that acts on
quantum input states with $\mathcal{O}(\log n)$ layers. Abbas et
al.~\cite{abbas2021} argued from an effective-dimension measure that some
quantum networks have higher capacity than comparable classical ones;
capacity alone, however, does not guarantee better accuracy on a real task.
Quantum kernel methods~\cite{havlicek2019} and data re-uploading
circuits~\cite{perezsalinas2020} are further alternatives, discussed as
future work.

\section{Trainability and Circuit Quality}
\label{sec:rw_quality}
McClean et al.~\cite{mcclean2018} showed that for sufficiently random
circuits the gradient variance vanishes exponentially with the number of
qubits, $\mathrm{Var}[\partial C/\partial\theta_j] \in \mathcal{O}(2^{-\alpha
n})$ --- a \emph{barren plateau}. Cerezo et al.~\cite{cerezo2021} showed that
global costs can plateau even for shallow circuits, whereas local costs
remain trainable. Sim et al.~\cite{sim2019} proposed two descriptors of a
circuit: \emph{expressibility}, the Kullback--Leibler divergence between the
distribution of fidelities $F$ of pairs of its random output states and that
of Haar-random states, $P_{\mathrm{Haar}}(F) = (N-1)(1-F)^{N-2}$ with $N =
2^n$ (lower is more expressive); and \emph{entangling capability}, the mean
Meyer--Wallach measure~\cite{meyer2002}
\begin{equation}
Q(\ket{\psi}) = 2\left(1 - \frac{1}{n}\sum_{k=1}^{n}\mathrm{Tr}\!\left(\rho_k^2\right)\right),
\end{equation}
where $\rho_k$ is the reduced state of qubit $k$ ($Q = 0$ for product
states). Holmes et al.~\cite{holmes2022} proved that the two properties
conflict: the closer a circuit's states are to Haar-random, the smaller its
gradients. On hardware, noise can additionally be reduced by error
mitigation such as zero-noise extrapolation~\cite{temme2017}.

\section{Research Gap}
Hybrid models for medical images are usually reported with a single run
and compared with a classical model that differs in size or structure, the
two hybrid placements are rarely compared on the same data, and the
circuits are seldom characterised beyond accuracy. This project evaluates
both placements on one TB dataset against matched controls, with repeated
seeds, several training-set sizes and paired significance tests, and
combines these results with an analysis of the circuits' expressibility,
trainability and noise robustness.
